\documentclass[12pt, a4paper]{article}

% ── Paketi ─────────────────────────────────────────────────────────────────────
\usepackage[utf8]{inputenc}
\usepackage[T1]{fontenc}
\usepackage{babel}
\usepackage{geometry}
\usepackage{graphicx}
\usepackage{hyperref}
\usepackage{booktabs}
\usepackage{array}
\usepackage{enumitem}
\usepackage{titlesec}
\usepackage{fancyhdr}
\usepackage{setspace}
\usepackage{parskip}
\usepackage{xcolor}
\usepackage{cite}
\usepackage{mdframed}

% ── Margine ────────────────────────────────────────────────────────────────────
\geometry{
  top=2.5cm, bottom=2.5cm,
  left=3cm,  right=2.5cm,
  headheight=15pt
}

% ── Prored ─────────────────────────────────────────────────────────────────────
\onehalfspacing

% ── Zaglavlje i podnožje ───────────────────────────────────────────────────────
\pagestyle{fancy}
\fancyhf{}
\rhead{\small Uticaj veštačke inteligencije na prirodu}
\lhead{\small Nikola Milović}
\cfoot{\thepage}
\renewcommand{\headrulewidth}{0.4pt}

% ── Boje ───────────────────────────────────────────────────────────────────────
\definecolor{forestgreen}{RGB}{44, 95, 45}
\definecolor{mossgreen}{RGB}{151, 188, 98}
\definecolor{lightmoss}{RGB}{212, 237, 170}
\definecolor{warningred}{RGB}{180, 50, 50}

% ── Naslovi sekcija ────────────────────────────────────────────────────────────
\titleformat{\section}
  {\large\bfseries\color{forestgreen}}
  {\thesection.}{0.5em}{}[\vspace{-0.3em}\rule{\linewidth}{0.4pt}]

\titleformat{\subsection}
  {\normalsize\bfseries\color{forestgreen!80!black}}
  {\thesubsection}{0.5em}{}

\titleformat{\subsubsection}
  {\normalsize\itshape\color{forestgreen!60!black}}
  {\thesubsubsection}{0.5em}{}

% ── Istaknuti okvir ────────────────────────────────────────────────────────────
\newmdenv[
  backgroundcolor=lightmoss!40,
  linecolor=forestgreen,
  linewidth=1pt,
  leftmargin=0pt,
  rightmargin=0pt,
  innerleftmargin=10pt,
  innerrightmargin=10pt,
  innertopmargin=8pt,
  innerbottommargin=8pt,
  skipabove=8pt,
  skipbelow=8pt
]{greenbox}

% ── Hiperlinks ─────────────────────────────────────────────────────────────────
\hypersetup{
  colorlinks=true,
  linkcolor=forestgreen,
  citecolor=forestgreen,
  urlcolor=forestgreen
}

% ═══════════════════════════════════════════════════════════════════════════════
\begin{document}

% ── Naslovna strana ────────────────────────────────────────────────────────────
\begin{titlepage}
  \centering
  \vspace*{2cm}

  {\Large \textbf{Univerzitet u Beogradu}\\[0.3em]
  \large Matematički fakultet}

  \vspace{1.5cm}

  \rule{\linewidth}{1pt}\\[0.4em]
  {\LARGE \textbf{Uticaj veštačke inteligencije na prirodu}}\\[0.3em]
  \rule{\linewidth}{1pt}

  \vspace{0.5cm}
  {\large \textit{Seminarski rad iz predmeta Računarstvo i društvo (RID)}}

  \vspace{2.5cm}

  \begin{tabular}{rl}
    \textbf{Student:}        & Nikola Milović \\[0.3em]
    \textbf{Broj indeksa:}   & mi21107 \\[0.3em]
    \textbf{Studijska god.:} & 2025/2026 \\[0.3em]
    \textbf{Mentor:}         & prof.\ dr Sana Stojanović \\
  \end{tabular}

  \vfill
  {\large Beograd, jun 2026.}
\end{titlepage}

\newpage
\tableofcontents
\newpage

% ═══════════════════════════════════════════════════════════════════════════════
\section{Uvod}

Veštačka inteligencija (VI) jedno je od najbrže rastućih tehnoloških polja
21.\ veka. Njena primena proširila se na gotovo sve oblasti ljudskog delanja
--- od medicine i obrazovanja, do transporta i zabave. Svakodnevno se pojavljuju
novi načini primene koji menjaju industrije i društvene navike, a tempo razvoja
čini se sve bržim. Međutim, jedno od najvažnijih i istovremeno
najkontroverznijih pitanja savremenog doba jeste: \textbf{kakav uticaj
veštačka inteligencija ima na prirodu i životnu sredinu?}

Ekološka kriza --- od klimatskih promena i gubitka biodiverziteta, do
zagađenja tla, vode i vazduha --- zahteva hitna i efikasna rešenja. Postavlja
se pitanje da li VI može biti deo tih rešenja, ili je sama deo problema.
Odgovor, kao što ćemo videti, nije jednoznačan.

Ovaj seminarski rad bavi se upravo tom temom. Cilj je da istražimo kako VI
može biti moćan saveznik u zaštiti planete, ali i kako njena nekontrolisana
upotreba može naneti ozbiljnu štetu ekosistemima. Rad sagledava obe strane
medalje: benefite koje VI pruža u oblastima poput monitoringa biodiverziteta,
pametne poljoprivrede i klimatskog modeliranja, ali i negativne posledice koje
prate razvoj VI sistema --- pre svega ogroman ugljenični otisak, potrošnju
vode i elektronski otpad.

Posebna pažnja posvećena je Jevons paradoksu --- fenomenu koji objašnjava
zašto tehnološka efikasnost ne vodi nužno do smanjenja ukupne potrošnje
resursa --- i etičkim dilemama oko kontrole podataka i greenwashinga. Na
kraju, razmatramo budućnost Zelene VI i mere koje su potrebne da bi VI zaista
postala saveznik, a ne neprijatelj planete.

% ═══════════════════════════════════════════════════════════════════════════════
\section{Veštačka inteligencija --- kratak pregled}

\subsection{Definicija i paradigme}

Veštačka inteligencija odnosi se na sposobnost računarskih sistema da
obavljaju zadatke koji tipično zahtevaju ljudsku inteligenciju: učenje,
rezonovanje, prepoznavanje obrazaca, donošenje odluka i razumevanje prirodnog
jezika. Termin je skovao John McCarthy 1956.\ godine, ali prava revolucija
nastupila je tek u poslednjoj deceniji, kada su dostupnost ogromnih skupova
podataka i napredni hardver omogućili dramatičan napredak.

Glavne paradigme savremene VI uključuju:

\begin{itemize}[leftmargin=1.5em]
  \item \textbf{Mašinsko učenje (Machine Learning)} --- algoritmi koji uče iz
        podataka bez eksplicitnog programiranja, pronalazeći obrasce u
        ogromnim skupovima informacija;
  \item \textbf{Duboko učenje (Deep Learning)} --- višeslojne neuronske mreže
        sposobne za prepoznavanje izuzetno kompleksnih obrazaca u slici, zvuku
        i tekstu;
  \item \textbf{Obrada prirodnog jezika (NLP)} --- razumevanje i generisanje
        ljudskog teksta, što je osnova velikih jezičkih modela;
  \item \textbf{Računarska vizija} --- interpretacija i analiza vizuelnih
        podataka, od prepoznavanja lica do detekcije bolesti na satelitskim
        snimcima.
\end{itemize}

\subsection{Veliki jezički modeli i ekološki otisak}

Ubrzani razvoj velikih jezičkih modela (LLM --- Large Language Models),
poput GPT i Claude serije, doneo je revoluciju u načinu na koji interagujemo
sa tehnologijom. Ovi modeli treniraju se na ogromnim skupovima podataka koji
mogu sadržati stotine milijardi reči, a sam proces treniranja zahteva mesece
rada na hiljadama specijalizovanih procesora \cite{strubell2019energy}.

Paradoks savremene VI leži u tome što isti sistemi koji mogu pomoći u zaštiti
planete istovremeno ostavljaju ogroman ekološki otisak. Razumevanje ovog
paradoksa ključno je za svaku ozbiljnu raspravu o ulozi VI u ekologiji.

% ═══════════════════════════════════════════════════════════════════════════════
\section{Pozitivni uticaji VI na prirodu}

\subsection{Praćenje i zaštita biodiverziteta}

Jedna od najznačajnijih primena VI u zaštiti prirode jeste monitoring
biodiverziteta. Tradicionalne metode praćenja životinjskih vrsta bile su
spore, skupe i ograničene geografskim dometom istraživača --- tim naučnika
mogao je pokriti svega nekoliko kvadratnih kilometara u jednoj ekspediciji.
Veštačka inteligencija menja ovu sliku radikalno.

\subsubsection{BirdNET --- prepoznavanje vrsta po zvuku}

Projekat \textit{BirdNET}, razvijen na Cornell Lab of Ornithology u saradnji
sa Chemnitz University of Technology, koristi duboko učenje za identifikaciju
ptičjih vrsta na osnovu zvuka \cite{kahl2021birdnet}. Sistem je do 2024.\
godine identifikovao više od 6.000 vrsta ptica u više od 190 zemalja i pomogao
naučnicima da mapiraju migracije i procene stanja populacija u realnom
vremenu. Ono što bi timu ornitologa oduzelo mesece terenskog rada, BirdNET
obavlja automatski analizom audio snimaka postavljenih senzora.

\subsubsection{Global Forest Watch --- satelitski nadzor šuma}

Organizacija \textit{Global Forest Watch}, uz pomoć VI algoritama, analizira
satelitske snimke i detektuje krčenje šuma gotovo u realnom vremenu ---
ažuriranja stižu svakih 48 sati. Ovo omogućuje vlastima i ekološkim
organizacijama da brzo reaguju na ilegalne seče i požare. Do 2024.\ godine,
sistem je pokrio praktično sve šumske površine planete, pružajući dragoceni
uvid u trend gubitka šuma koji inače ne bi bio dostupan decenijama.

\subsubsection{Zaštita ugroženih vrsta i borba protiv krivolova}

VI modeli koriste se za predviđanje kretanja divljih životinja i identifikaciju
krivolova. Sistem PAWS (\textit{Protection Assistant for Wildlife Security})
analizira istorijske podatke o krivolovu zajedno sa klimatskim uslovima i
karakteristikama terena kako bi predvideo verovatne lokacije budućih
incidenata. Čuvari parkova dobijaju preporuke gde da patroliraju, drastično
povećavajući efikasnost zaštite. Slično, Microsoft AI for Earth finansirao je
projekte koji su pomogli u identifikaciji više od 30.000 kitova u okeanima,
koristeći VI analizu hidroakustičkih snimaka \cite{norouzzadeh2018automatically}.

\subsection{Borba protiv klimatskih promena}

Klimatske promene predstavljaju jednu od najvećih egzistencijalnih pretnji po
planet. VI igra sve važniju ulogu u razumevanju i ublažavanju ove pretnje
\cite{rolnick2019tackling}.

\textbf{Klimatsko modeliranje:} Moderni klimatski modeli oslanjaju se na VI
da procesiraju ogromne količine podataka sa meteoroloških stanica, satelita i
okeanskih bova. Ovakvi modeli su daleko precizniji od prethodnih generacija i
omogućuju bolje predviđanje ekstremnih vremenskih pojava. Google Flood Hub
koristi VI za predviđanje poplava do sedam dana unapred i već pokriva 80
zemalja, direktno upozoravajući lokalno stanovništvo.

\textbf{Optimizacija energetskih sistema:} DeepMind (Google) primenio je VI
za optimizaciju sistema za hlađenje u serverskim centrima, postižući smanjenje
energetske potrošnje za hlađenje do \textbf{40\%}. Isti princip se primenjuje
za upravljanje vetroelektranama --- VI predviđa produkciju vetra 36 sati
unapred, povećavajući vrednost energije za oko 20\%.

\subsection{Pametna poljoprivreda}

Poljoprivreda je jedan od najvećih potrošača vode i jednih od najvećih
emitera gasova staklene bašte. VI transformiše ovaj sektor kroz preciznu
poljoprivredu na više načina:

\begin{itemize}[leftmargin=1.5em]
  \item \textbf{Precizno navodnjavanje:} Senzori u kombinaciji sa VI
        algoritmima mere vlažnost zemljišta i vremensku prognozu, dozirajući
        vodu tačno tamo gde je potrebna. Studije pokazuju da ovakvi sistemi
        smanjuju potrošnju vode do \textbf{50\%} u poređenju sa tradicionalnim
        metodama navodnjavanja;
  \item \textbf{Ciljana primena pesticida:} John Deere See \& Spray koristi
        računarsku viziju da identifikuje korove i primeni herbicide samo na
        zaražena mesta --- smanjujući upotrebu hemikalija za \textbf{66\%};
  \item \textbf{Dijagnostika bolesti useva:} Aplikacija Plantix koristi VI
        za prepoznavanje bolesti biljaka na osnovu fotografija, pomažući
        farmerima u 70+ zemalja i sa više od 10 miliona korisnika.
\end{itemize}

\subsection{Upravljanje otpadom i reciklažom}

Roboti opremljeni VI sistemima računarske vizije mogu sortirati reciklabilni
materijal znatno brže i tačnije od ljudskih radnika. Kompanija AMP Robotics
razvila je sistem koji sortira do 80 predmeta u minuti, što je daleko iznad
ljudskih mogućnosti. Pored toga, VI analizira podatke o potrošnji u
maloprodajnim lancima i restoranima, pomažući u predviđanju tražnje i
smanjenju bacanja hrane --- što ima značajan ekološki i ekonomski efekat.

% ═══════════════════════════════════════════════════════════════════════════════
\section{Negativni uticaji VI na prirodu}

\subsection{Energetska potrošnja i ugljenični otisak}

Razvoj i eksploatacija VI sistema zahtevaju ogromne količine energije. Ovo je
možda najznačajniji i najčešće zanemarivan negativni uticaj VI na prirodu.

Istraživanje iz 2019.\ godine (Strubell et al.) pokazalo je da treniranje
jednog velikog NLP modela može emitovati toliko CO\textsubscript{2} koliko
\textbf{pet automobila u toku celog veka upotrebe} \cite{strubell2019energy}.
Noviji modeli poput GPT-4 su daleko veći i zahtevaju još više resursa. Prema
procenama, globalni serverski centri troše između \textbf{200 i 250
teravat-sati} električne energije godišnje --- što odgovara godišnjoj potrošnji
pojedinih zemalja srednje veličine \cite{patterson2021carbon}.

Studija objavljena 2024.\ u Nature časopisu analizirala je podatke iz 67
zemalja i zaključila da, iako VI može smanjiti ekološki otisak kroz
optimizaciju, njen sopstveni ugljenični trag raste eksponencijalnom brzinom
\cite{nature2024ecological}.

\subsection{Elektronski otpad i retki metali}

Hardver neophodan za VI --- GPU čipovi, memorijski moduli, specijalizovani
procesori (TPU, NPU) --- ima relativno kratak životni vek i teško se reciklira.

\textbf{Rudarenje retkih metala:} Proizvodnja naprednih čipova zahteva čak 17
različitih retkih metala, uključujući kobalt, litijum, tantal i neodim.
Rudarenje ovih materijala uzrokuje devastaciju ekosistema --- na primer, oko
70\% svetskog kobalta potiče iz Demokratske Republike Kongo, gde rudnici
nanose ozbiljne ekološke i socijalne štete.

\textbf{Brza zastarelost:} Brz tempo razvoja VI hardvera znači da oprema
zastareva u roku od svega 3 do 5 godina, doprinoseći rastućoj krizi
elektronskog otpada. Procenjuje se da globalna količina e-otpada premašuje
\textbf{60 miliona tona} godišnje --- više nego što iznosi težina Velikog
kineskog zida.

\subsection{Potrošnja vode}

Serverski centri troše ogromne količine vode za hlađenje --- problem koji je
dugo ostajao u senci energetskih pitanja. Istraživanje Li et al.\ (2023)
pokazuje da trening jednog velikog jezičkog modela može potrošiti i do
\textbf{700.000 litara} vode \cite{li2023thirsty}. Na operativnom nivou,
procenjuje se da korisnik potroši oko 500 ml vode za svakih 20--50 upita
prosečnom VI asistentumu. Ovo je posebno problematično za serverske centre
locirane u regionima koji se suočavaju sa nestašicom vode.

\subsection{Jevons paradoks --- kada efikasnost povećava potrošnju}

Jedan od najvažnijih, a često zanemarenih fenomena u diskusiji o VI i ekologiji
jeste \textbf{Jevons paradoks}. William Stanley Jevons primetio je 1865.\
godine da povećanje efikasnosti korišćenja resursa paradoksalno dovodi do
povećanja ukupne potrošnje tog resursa.

\begin{greenbox}
\textit{,,Povećanje efikasnosti korišćenja resursa paradoksalno dovodi do
povećanja ukupne potrošnje tog resursa.''} \\
\hfill --- William Stanley Jevons, 1865.
\end{greenbox}

U kontekstu VI, ovaj paradoks se manifestuje na više načina. Efikasniji VI
modeli za optimizaciju transporta ne smanjuju ukupan broj prevoženih
kilometara --- naprotiv, čine prevoz pristupačnijim, što povećava tražnju.
Jeftiniji letovi zahvaljujući VI optimizaciji cena dovode do više putnika i
više letova, a time i do većih ukupnih emisija. Sami VI modeli --- kako
postaju efikasniji --- primenjuju se u sve više oblasti, pa ukupna potrošnja
energije raste uprkos manjoj potrošnji po jednom upitu.

% ═══════════════════════════════════════════════════════════════════════════════
\section{Poređenje VI i tradicionalnih pristupa}

Da bismo bolje razumeli vrednost VI u ekologiji, korisno je direktno uporediti
efikasnost VI rešenja sa tradicionalnim metodama u ključnim oblastima:

\begin{table}[h!]
\centering
\caption{Poređenje VI i tradicionalnih metoda u zaštiti životne sredine}
\renewcommand{\arraystretch}{1.3}
\begin{tabular}{>{\bfseries\color{forestgreen}}p{2.8cm} p{4cm} p{5cm}}
\toprule
Oblast & Tradicionalno & Sa VI \\
\midrule
Monitoring šuma     & Terenske ekspedicije (meseci) & Sateliti + VI, ažuriranje svaka 48h \\
Identifikacija vrsta & Ekspert na terenu, ograničen doseg & BirdNET: 6.000+ vrsta automatski \\
Navodnjavanje       & Fiksni vremenski rasporedi & Senzori + VI, ušteda do 50\% vode \\
Prognoza poplava    & 1--2 dana unapred & 7 dana unapred, 80 zemalja \\
Sortiranje otpada   & $\sim$30 predmeta/min (čovek) & $\sim$80 predmeta/min (robot) \\
Otkrivanje krivolova & Nasumične patrole & PAWS: prediktivno patroliranje \\
\bottomrule
\end{tabular}
\end{table}

Tabela jasno pokazuje da VI u gotovo svim oblasticama donosi merljivo poboljšanje
efikasnosti. Ipak, važno je naglasiti da ova poboljšanja dolaze uz ekološke
troškove koji moraju biti uzeti u obzir pri donošenju odluka o implementaciji.

% ═══════════════════════════════════════════════════════════════════════════════
\section{Etička razmatranja}

Primena VI u zaštiti prirode nameće brojna etička pitanja koja se ne smeju
zanemariti \cite{vinuesa2020role}.

\subsection{Ko kontroliše podatke o prirodi?}

Sakupljanje podataka o biodiverzitetu i ekosistemima često se odvija u
zemljama Globalnog juga --- u Africi, Aziji i Latinskoj Americi, gde se
nalazi najveći deo svetskog biodiverziteta. Međutim, koristi od VI analize
ovih podataka uglavnom ubiru korporacije i institucije Globalnog severa.
Ovo stvara neravnopravan odnos koji mora biti adresiran kroz pravedne modele
deljenja prihoda i uključivanje lokalnih zajednica u donošenje odluka.

\subsection{Greenwashing --- zeleni imidž nasuprot stvarnom uticaju}

Postoji rastuća zabrinutost da neke kompanije naglašavaju ekološke primene
VI kako bi odvratile pažnju od ogromnog ekološkog otiska njihove
infrastrukture. Ovaj fenomen --- poznat kao \textit{greenwashing} --- posebno
je opasan jer može usporiti donošenje regulatornih mera. Tipičan primer je
situacija u kojoj kompanija objavljuje projekat za zaštitu koraljnih grebena
koristeći VI, dok istovremeno njeni serverski centri emituju milione tona
CO\textsubscript{2} godišnje.

\subsection{Autonomne odluke i odgovornost}

Kada VI sistemi donose autonomne odluke o upravljanju ekosistemima ---
na primer, o doziranju hemikalija u preciznoj poljoprivredi ili o lokacijama
patroliranja u zaštićenim oblastima --- postavlja se pitanje odgovornosti u
slučaju neželjenih posledica. Ko snosi odgovornost ako VI sistem pogrešno
klasifikuje vrstu i preporuči njeno uklanjanje? Ova pitanja još uvek nisu
pravno razrešena u većini zemalja.

\subsection{Nadzor divljih životinja i privatnost}

Intenzivno praćenje divljih životinja pomoću VI sistema može poremetiti
njihovo prirodno ponašanje. Postavljanje senzora, kamera i dronova u
ekosisteme nije bezopasno, a etička granica između korisnog monitoringa i
invazivnog nadzora nije uvek jasna. Pored toga, podaci prikupljeni o
kretanju životinja --- naročito ugroženih vrsta --- mogu biti zloupotrebljeni
od strane krivolovaca ako ne budu adekvatno zaštićeni.

% ═══════════════════════════════════════════════════════════════════════════════
\section{Budući pravci --- put ka Zelenoj VI}

Da bi VI zaista postala snaga za dobro u zaštiti planete, potrebni su koraci
na više fronova. Evo ključnih pravaca za narednu deceniju:

\subsection{Zelena VI (Green AI)}

Akademska zajednica i industrija moraju staviti veći naglasak na energetsku
efikasnost algoritama --- ne samo na njihovu tačnost i performanse. Istraživanja
u oblasti \textit{neuromorphic computing} --- računarskih sistema koji
oponašaju energetsku efikasnost ljudskog mozga --- i kvantnih računara
obećavaju dramatično smanjenje energetske potrošnje. Već sada postoji trend
ka manjim, efikasnijim modelima koji postižu slične rezultate uz daleko manji
ekološki trag.

\subsection{Obnovljivi izvori energije}

Serverski centri treba da budu napajani isključivo obnovljivim izvorima
energije --- solarnom, vetro i hidroenergijom. Neke kompanije već idu tim
putem: Google je od 2017.\ odgovarao za 100\% potrošnje energije kupovinom
obnovljive energije, a Microsoft se obavezao na ugljičnu negativnost do 2030.\
Međutim, tempo ove tranzicije u celoj industriji ostaje nedovoljan.

\subsection{Regulativa i transparentnost}

Potreban je zakonodavni okvir koji će propisati obavezno izveštavanje o
ugljeničnom otisku VI sistema, slično kao što se to čini za emisije u
industriji. Evropski AI Act iz 2023.\ sadrži neke odredbe u tom smeru, ali
su one još uvek nedovoljne. Potrebni su globalni standardi koji obuhvataju
ceo životni ciklus VI sistema --- od rudarenja sirovina do zbrinjavanja
hardvera.

\subsection{Međunarodna saradnja}

Ekološki problemi ne poznaju granice, pa ni rešenja ne smeju biti ograničena
na nacionalne okvire. Potrebni su globalni sporazumi o standardima za
„ekološku VI" --- sporazumi koji uključuju i razvijene i zemlje u razvoju,
i koji obezbeđuju pravičnu raspodelu koristi i tereta zelene tranzicije.

% ═══════════════════════════════════════════════════════════════════════════════
\section{Zaključak}

Veštačka inteligencija stoji pred nama kao jedna od najznačajnijih tehnologija
u istoriji čovečanstva --- sa potencijalom da pomogne u rešavanju nekih od
najtežih ekoloških izazova sa kojima se suočavamo. Od praćenja ugroženih vrsta
i predviđanja klimatskih katastrofa, do optimizacije energetskih sistema i
precizne poljoprivrede --- primene su raznovrsne, impresivne i sve dostupnije.

Međutim, ova ista tehnologija dolazi sa sopstvenim, često potcenjenim ekološkim
cehom. Ugljenični otisak treniranja modela, potrošnja vode u serverskim
centrima, elektronski otpad i devastacija ekosistema u rudnicima retkih metala
--- sve su to stvarni, merljivi problemi koji zahtevaju sistemska rešenja.
Jevons paradoks nas upozorava da tehnološka efikasnost sama po sebi nije
dovoljna --- neophodni su svesne politike i regulativa.

\begin{greenbox}
Odgovor ne leži u odbacivanju VI, već u \textbf{svesnom i odgovornom razvoju}
ove tehnologije. Budućnost mora biti ona u kojoj VI sistemi troše minimalnu
energiju, rade na obnovljivim izvorima, i čiji je pun životni ciklus ---
od rudarenja sirovina do zbrinjavanja otpada --- ekološki odgovoran.
\end{greenbox}

Planeta ne može čekati. Veštačka inteligencija može biti deo rešenja ---
ali samo ako sebi kao društvo postavimo visoke standarde i budemo spremni
da ih sprovedemo u delo.

% ═══════════════════════════════════════════════════════════════════════════════
\newpage
\bibliographystyle{plain}
\bibliography{literatura}

\end{document}
